\subsection{\texorpdfstring{Planitud y propiedad $(S_k)$.}{Planitud y propiedad (S\string_k)}}

\begin{proposition}[6.4.1]
\label{IV.2.s6.4.1-fr}
Sean $X$, $Y$ dos preesquemas localmente noetherianos, $f:X\to Y$ un
morfismo, $\mathscr E$ un $\mathcal O_Y$-módulo coherente y
$\mathscr F$ un $\mathcal O_X$-módulo coherente y $f$-plano.
\begin{enumerate}
\item[\textup{(i)}] Sea $x\in\operatorname{Supp}(\mathscr F)$ tal que
$\mathscr E\otimes_Y\mathscr F$ posea la propiedad $(S_k)$ en el punto $x$;
entonces $\mathscr E$ posee la propiedad $(S_k)$ en el punto $y=f(x)$.
\item[\textup{(ii)}] Supongamos que, para todo
$y\in f(\operatorname{Supp}(\mathscr F))$,
$\mathscr F_y=\mathscr F\otimes_{\mathcal O_Y}\boldsymbol{k}(y)$ posee la propiedad
$(S_k)$; entonces, si para un punto
$y\in f(\operatorname{Supp}(\mathscr F))$, $\mathscr E$ posee la propiedad
$(S_k)$ en el punto $y$, $\mathscr E\otimes_Y\mathscr F$ posee la propiedad
$(S_k)$ en todo punto de $f^{-1}(y)$.
\end{enumerate}
\end{proposition}

(i) Se sabe $(\mathbf{Err}_{\mathrm{III}},30)$ que existe un subpreesquema
cerrado $X'$ de $X$ cuyo espacio subyacente es
$\operatorname{Supp}(\mathscr F)$ y un $\mathcal O_{X'}$-módulo coherente
$\mathscr F'$ tales que $\mathscr F=j_*(\mathscr F')$, donde $j$ es la
inyección canónica $X'\to X$; se puede sustituir $X$ por $X'$ y $\mathscr F$
por $\mathscr F'$, es decir, suponer que
$\operatorname{Supp}(\mathscr F)=X$. Sea entonces $y'$ una generalización de
$y$ en $Y$; se sabe (2.3.4) que existe una generalización $x'$ de $x$ en $X$
tal que $y'=f(x')$; se puede, además, suponer que $x'$ es punto genérico de
una componente irreducible de $f^{-1}(y')$; se tiene, por tanto,
$\dim(\mathscr F_{x'}\otimes_{\mathcal O_{y'}}\boldsymbol{k}(y'))=0$. En virtud de (6.1.2),
se tiene, poniendo $\mathscr G=\mathscr E\otimes_Y\mathscr F$,
\[
\dim(\mathscr G_{x'})=\dim(\mathscr E_{y'})
\]
y, en virtud de (6.3.1),
\[
\operatorname{prof}(\mathscr G_{x'})=\operatorname{prof}(\mathscr E_{y'})
\]
teniendo en cuenta que la profundidad es siempre a lo sumo igual a la
dimensión. Por hipótesis, se tiene
\[
\operatorname{prof}(\mathscr G_{x'})
\geq\inf(k,\dim(\mathscr G_{x'}))
\]
luego
\[
\operatorname{prof}(\mathscr E_{y'})
\geq\inf(k,\dim(\mathscr E_{y'}))
\]
lo que demuestra la primera afirmación.

\oldpage[IV-2]{142}

(ii) Como para toda generalización $x'$ de $x\in f^{-1}(y)$,
$y'=f(x')$ es una generalización de $y$, es posible limitarse a verificar que,
si $x\in\operatorname{Supp}(\mathscr G)$ y $f(x)=y$, se tiene
$\operatorname{prof}(\mathscr G_x)\geq\inf(k,\dim(\mathscr G_x))$; resulta de
(6.1.2) y de (6.3.1) que se tiene
\[
\dim(\mathscr G_x)=\dim(\mathscr E_y)+
\dim(\mathscr F_x\otimes_{\mathcal O_y}\boldsymbol{k}(y))
\]
\[
\operatorname{prof}(\mathscr G_x)=\operatorname{prof}(\mathscr E_y)+
\operatorname{prof}(\mathscr F_x\otimes_{\mathcal O_y}\boldsymbol{k}(y)).
\]

Por hipótesis, se tiene
\[
\operatorname{prof}(\mathscr E_y)\geq\inf(k,\dim(\mathscr E_y))
\]
\[
\operatorname{prof}(\mathscr F_{\boldsymbol{x}}\otimes_{\mathcal O_y}\boldsymbol{k}(y))
\geq\inf(k,\dim(\mathscr F_x\otimes_{\mathcal O_y}\boldsymbol{k}(y)))
\]
de donde, sumando miembro a miembro,
\begin{align*}
\operatorname{prof}(\mathscr G_{\boldsymbol{x}})
&\geq\inf(k,\dim(\mathscr E_y))
+\inf(k,\dim(\mathscr F_{\boldsymbol{x}}\otimes_{\mathcal O_y}\boldsymbol{k}(y)))\geq{}\\
&\geq\inf\bigl(k,\dim(\mathscr E_y)
+\dim(\mathscr F_x\otimes_{\mathcal O_y}\boldsymbol{k}(y))\bigr)
=\inf(k,\dim(\mathscr G_{\boldsymbol{x}}))
\end{align*}
lo que demuestra (ii).
